% problem-type: nhiet-mien-co-bien
% order: 60
% Bo cuc: De vi du -> Cong thuc -> De + Bai giai

\section{Nhiệt trên miền có biên --- thác triển + Poisson}

% ---------------------------------------------------------------
\subsection*{Đề bài ví dụ}
% ---------------------------------------------------------------
Giải $u_t=4\Delta u$ trên dải $0<x<\pi,\ y>0$, $\;u(x,y,0)=\sin^3(x)\chi_{[0,1]}(y)$,
với biên $u(0,y,t)=u(\pi,y,t)=u(x,0,t)=0$.

\emph{\small(Nguồn: Đề thi MAT3365 --- PTĐHR, HK2 2023--2024, K67 Toán tin, Câu 2. Lời giải dưới đây tự soạn.)}

% ---------------------------------------------------------------
\subsection*{Nhận ra dạng này khi nào?}
% ---------------------------------------------------------------
\begin{itemize}
  \item Nhiệt trên miền \textbf{có biên} (đoạn, dải, nửa mặt phẳng, hộp), biên Dirichlet/Neumann.
  \item Ý tưởng: \textbf{thác triển} dữ liệu ra toàn $\mathbb{R}^n$ để tự thỏa biên, rồi dùng Poisson.
\end{itemize}
\begin{meobox}
Dirichlet $u=0\Rightarrow$ thác triển lẻ; Neumann $u_x=0\Rightarrow$ chẵn; đoạn $\Rightarrow$ tuần hoàn.
\end{meobox}

% ---------------------------------------------------------------
\subsection*{Công thức \& phương pháp}
% ---------------------------------------------------------------
\subsubsection*{Thác triển khử biên}
% technique: thac-trien
\textbf{Thác triển để khử điều kiện biên} (nhiệt/sóng trên nửa không gian hay đoạn).
Mở rộng dữ liệu ra toàn $\mathbb{R}^n$ sao cho \emph{tự} thỏa điều kiện biên, rồi áp công thức Poisson/d'Alembert:
\begin{itemize}
  \item Biên Dirichlet $u=0$ tại $x=0$ $\Rightarrow$ thác triển \textbf{lẻ} qua $x=0$.
  \item Biên Neumann $u_x=0$ tại $x=0$ $\Rightarrow$ thác triển \textbf{chẵn} qua $x=0$.
  \item Miền là đoạn/hộp $\Rightarrow$ thác triển \textbf{tuần hoàn} (lẻ/chẵn tùy biên).
\end{itemize}
Ví dụ: $u=0$ tại $x=0$ và $x=\pi$ $\Rightarrow$ thác triển lẻ, tuần hoàn chu kỳ $2\pi$.

\subsubsection*{Rồi dùng công thức Poisson}
% method: cong-thuc-poisson-nhiet
\begin{dieukienbox}
Dùng khi: phương trình \textbf{nhiệt} $u_t=k\Delta u$ trên toàn $\mathbb{R}^n$, dữ liệu $u(\cdot,0)=g$.
\end{dieukienbox}

\begin{nhobox}
\textbf{Công thức Poisson} (tích chập với nhân nhiệt $\Phi$):
\[
  u(x,t)=\int_{\mathbb{R}^n}\Phi(x-y,t)\,g(y)\,dy,\qquad
  \Phi(x,t)=\frac{1}{(4\pi k t)^{n/2}}\,e^{-|x|^2/(4kt)}.
\]
1D: $\displaystyle u(x,t)=\frac{1}{\sqrt{4\pi k t}}\int_{-\infty}^{\infty} e^{-(x-y)^2/(4kt)}\,g(y)\,dy.$
\end{nhobox}

\begin{meobox}
$g$ là hàm đặc trưng $\chi_{[a,b]}$ $\Rightarrow$ đổi biến $w=\frac{y-x}{\sqrt{4kt}}$, tích phân ra
hàm $\operatorname{erf}$ (xem Phụ lục B). Nhiều chiều: nhân nhiệt tách tích theo từng biến.
\end{meobox}
\emph{(Tính duy nhất nghiệm: dùng tích phân năng lượng --- Phụ lục C.)}

% ---------------------------------------------------------------
\subsection*{Đề bài \& bài giải}
% ---------------------------------------------------------------
\paragraph{Bước 1 --- thác triển khử biên.} Dữ liệu tách được $g=\sin^3(x)\cdot\chi_{[0,1]}(y)$, nên giải
tách theo $x$ và $y$. Theo $x$: $u=0$ tại $x=0,\pi\Rightarrow$ thác triển \textbf{lẻ}, tuần hoàn $2\pi$.
Theo $y$: $u=0$ tại $y=0\Rightarrow$ thác triển \textbf{lẻ} ra $y<0$: $\chi_{[0,1]}\to\chi_{[0,1]}-\chi_{[-1,0]}$.

\paragraph{Bước 2 --- phần $x$ (chuỗi sin, mỗi mốt tắt riêng).} Hạ bậc: $\sin^3 x=\tfrac34\sin x-\tfrac14\sin 3x$.
Mỗi $\sin(mx)$ là hàm riêng, qua nhiệt $u_t=4u_{xx}$ tắt theo $e^{-4m^2 t}$ ($m=1\to e^{-4t}$, $m=3\to e^{-36t}$):
\[
  X(x,t)=\tfrac34 e^{-4t}\sin x-\tfrac14 e^{-36t}\sin 3x=\frac{3e^{-4t}\sin x-e^{-36t}\sin 3x}{4}.
\]
(Kiểm tra $t\to0$: $X\to\tfrac{3\sin x-\sin3x}{4}=\sin^3 x$. Khớp.)

\paragraph{Bước 3 --- phần $y$ (Gauss $\to$ erf).} Poisson với $k=4$ (nhân $\tfrac{1}{\sqrt{16\pi t}}e^{-(y-Y)^2/16t}$),
dữ liệu lẻ $\chi_{[0,1]}-\chi_{[-1,0]}$ tách thành hai tích phân:
\[
  Y(y,t)=\Big[\int_0^1-\int_{-1}^0\Big]\tfrac{1}{\sqrt{16\pi t}}e^{-(y-Y)^2/16t}\,dY.
\]
Đổi biến $z=\tfrac{Y-y}{4\sqrt t}$: mỗi $\int_p^q=\tfrac12\big(\operatorname{erf}\tfrac{q-y}{4\sqrt t}-\operatorname{erf}\tfrac{p-y}{4\sqrt t}\big)$.
Dùng $\operatorname{erf}$ lẻ, gộp lại:
\[
  Y(y,t)=\operatorname{erf}\tfrac{y}{4\sqrt t}-\tfrac12\Big(\operatorname{erf}\tfrac{y+1}{4\sqrt t}+\operatorname{erf}\tfrac{y-1}{4\sqrt t}\Big).
\]
(Kiểm tra $t\to0$, $y=\tfrac12\in(0,1)$: $1-\tfrac12(1-1)=1$. Khớp $\chi_{[0,1]}$.)

\paragraph{Bước 4 --- nhân hai phần.}\leavevmode

\begin{nhobox}
\[
  \begin{aligned}
  u(x,y,t)=\ &\frac{3e^{-4t}\sin x-e^{-36t}\sin 3x}{4}\\[2pt]
  &\times\ \Big[\operatorname{erf}\tfrac{y}{4\sqrt t}-\tfrac12\big(\operatorname{erf}\tfrac{y+1}{4\sqrt t}+\operatorname{erf}\tfrac{y-1}{4\sqrt t}\big)\Big].
  \end{aligned}
\]
\end{nhobox}